\begin{helper}[A maximum balanced candidate has Tashkinov order]
Let $N$ be normal globally maximum data for $(T,e_0,\varphi)$, and let
$S$ be a \noderef{ScheideBalancedCandidate} for $N$ whose order is maximum
among all such candidates.  Then
\[
 |V(S)|=\noderef{TashkinovOrder}(G).
\]
\end{helper}
\begin{proof}
Write $p=N.\operatorname{height}$, $i=|V(S)|$, and
$t=\noderef{TashkinovOrder}(G)$.  By
\noderef{ScheideBalancedCandidate}, the hypothesis on $S$ supplies its
Tashkinov property, odd order, exact-trunk initial segment, used-color
restriction, and post-trunk equal-color pairing.  The global field of $N$,
exposed by \noderef{NormalGlobalMaximumTriple} and interpreted by
\noderef{GlobalMaximumTriple}, supplies both the fixed
\noderef{KTriple} $(G,e_0,\varphi)$ and $|V(T)|=t$.  In particular
$\chi'(G)=k+1$, so $k=\chi'(G)-1$.  Since its candidate Tashkinov field is
a \noderef{TashkinovTree} for $S$, its cardinality occurs in the finite
maximum defining $t$; hence $i\le t$.  We prove that strict inequality is
impossible.

Assume $i<t=|V(T)|$.  Choose the first displayed vertex $y_r$ of $T$ not
in $V(S)$.  The first edge identity of both sequences is $e_0$, and the
attachment equation in \noderef{TreeSequence}, together with looplessness,
says that their first two displayed vertices are the two distinct ends of
$e_0$.  Thus both belong to $V(S)$, so $r\ge2$, and every $y_a$ with
$a<r$ lies in $V(S)$.  The Tashkinov property of $T$ supplies a color
$\gamma$ and an index $a<r$ such that the exact original edge identity
$e_r$ has color $\gamma$ and $y_a$ misses $\gamma$.  The attachment
equation supplies $b<r$ such that the ends of $e_r$ are $y_r$ and $y_b$.

Construct a tree-sequence $S^+$, rather than suppressing the append.  By
the injective vertex map in \noderef{TreeSequence} and the definition
\noderef{TreeSequenceVertexSet}, $i=S.\operatorname{length}+1$, so $S$ has
length $i-1$.  Give $S^+$ length $i$, retain the old vertex and edge
functions on their old index ranges, set its last vertex equal to $y_r$,
and set its last edge identity equal to $e_r$.  The new vertex is fresh by
the choice of $r$.  Every old edge of $S$ has both endpoints in $V(S)$ by
\noderef{TreeSequence}, whereas $e_r$ has the endpoint $y_r$ outside
$V(S)$, so $e_r$ is a fresh edge identity.  Its other endpoint $y_b$ is an
old vertex, which proves the last attachment equation; all earlier
attachment equations and both injectivity clauses are inherited.  The old
Tashkinov witnesses remain valid, and $y_a$ is the required earlier
missing-color witness for the new edge.  Hence $S^+$ is a Tashkinov tree.
Its constructed vertex map is injective, so
\noderef{TreeSequenceVertexSet} gives it the even order $i+1$.  Also
$\gamma\in\noderef{TreeUsedColors}(\varphi,T)$ because $e_r$ is an
original edge at one-based position $r\ge2$.

Apply \noderef{TashTree} to the unchanged $k$-triple and $S^+$.  Its
vertex set $U$ is elementary in the sense of
\noderef{ColoredVertexSetElementary}.  The vertex $y_a$ misses $\gamma$,
and elementariness says that no second vertex of $U$ does.  Thus exactly
one vertex of $U$ misses $\gamma$.  By \noderef{MissingColors} and
\noderef{PresentColors}, every other vertex of $U$ is incident with at
least one $\gamma$-colored edge identity.  It is incident with at most one:
two distinct such identities would be adjacent by
\noderef{AdjacentEdges}, contradicting the proper-coloring clause of the
$k$-triple and \noderef{IsProperEdgeColoring}.  The number of incidence
pairs
\[
 (v,f),\qquad v\in U,\quad \noderef{Incident}(G,v,f),\quad
 \varphi(f)=\gamma,
\]
is therefore $|U|-1=i$, an odd number.

This count retains individual multigraph edge identities.  Partition the
$\gamma$-colored identities meeting $U$ into those with both endpoints in
$U$ and those with exactly one endpoint in $U$.  Say that the two classes
have $a_0$ and $b_0$ identities.  An internal identity contributes its two
distinct endpoints, by the loopless law in
\noderef{LooplessMultigraph}; a second-class identity contributes its
unique endpoint in $U$; all other identities contribute zero.  Hence
$i=2a_0+b_0$.  The second class is exactly the set of $\gamma$-colored
identities in \noderef{BoundaryEdges}$(G,U)$.  Since $i$ is odd, $b_0$ is
odd and therefore nonzero.  Choose one exact identity $f$ from this class,
with endpoints $u\in U$ and $w\notin U$.

Now explicitly construct $S^{++}$ by giving it length $i+1$, retaining all
entries of $S^+$, and placing $f$ and $w$ in its final edge and vertex
positions.  The vertex $w$ is fresh.  Every old edge identity has both ends
in $U$, whereas $f$ is a boundary identity, so $f$ is fresh.  Its endpoint
$u$ proves the new attachment equation.  All old Tashkinov witnesses are
unchanged, and the unique vertex of $U$ missing $\gamma$ is an earlier
witness for the new $\gamma$-colored edge $f$.  Thus $S^{++}$ is again a
Tashkinov tree-sequence.

It remains to check every field of
\noderef{ScheideBalancedCandidate}.  The constructed $S^{++}$ is Tashkinov
by the preceding paragraph.  By the exact-entry definitions in
\noderef{NormalTripleTrunk} and \noderef{TreeSequenceInitialSegment},
appending entries preserves the exact initial trunk.  All old used colors
obey the restriction for $S$, and both new edges have the same color
$\gamma\in\noderef{TreeUsedColors}(\varphi,T)$.  The exact-trunk condition
from \noderef{NormalTripleTrunk} and
\noderef{TreeSequenceInitialSegment} gives $p\le i$: the trunk has length
$p-1$, while $S$ has length $i-1$.  The odd-order candidate field says that
$i$ is odd.  Thus the two new edge positions are precisely the next one-based pair
$(e_i,e_{i+1})$: equivalently their zero-based positions are $i-1$ and
$i$, corresponding to the unique new even number $2j=i+1$ in the pairing
range $p<2j<i+2$.  Therefore the old equal-color pairs plus
$\varphi(e_r)=\varphi(f)=\gamma$ prove the entire pair condition.  Finally,
the constructed vertex map is injective, so
\noderef{TreeSequenceVertexSet} gives $|V(S^{++})|=i+2$, which is odd.
Thus all five clauses of \noderef{ScheideBalancedCandidate} hold for
$S^{++}$.  It is a balanced candidate of order $i+2$, contradicting the
maximality of $S$.  We conclude
$i\not<t$, and together with $i\le t$ this gives $i=t$.
\end{proof}
